% Chapter 7: Lie Groups
% Lee, \emph{Introduction to Smooth Manifolds} (2nd ed., GTM 218).
% Generated from the section extraction; nodes carry only Lean that is
% verified sorry-free. See ../../UPSTREAM_LEAN_AUDIT.md.

\begin{definition}
\label{def:7-46-extra-1}
\lean{LieGroup}

\mathlibok
A Lie group is a boundaryless smooth manifold $G$ with a group structure for which multiplication $m:G\times G\to G$, $m(g,h)=gh$, and inversion $i:G\to G$, $i(g)=g^{-1}$, are smooth. Hence every Lie group is a topological group.
\end{definition}

\begin{notation}
\label{not:7-46-extra-2}
\lean{LieAddGroup}

\mathlibok
Products in a Lie group are written multiplicatively, with identity $e$;
additive notation and $0$ are used for abelian examples such as $\mathbb R^n$,
and $I_n$ denotes the identity matrix.
\end{notation}

\begin{proposition}
\label{prop:7-1}
\lean{lieGroup_of_contMDiff_mul_inv}
\leanok


\dcref{ch7:7.1}
If a smooth manifold $G$ has a group structure and $(g,h)\mapsto gh^{-1}$ is smooth, then $G$ is a Lie group.
\end{proposition}


\begin{example}[Lie Groups]
\label{ex:7-3}
\lean{ContinuousLinearEquiv.unitsEquiv}

\mathlibok
\dcref{ch7:7.3}
The following are Lie groups.

(a) $\mathrm{GL}(n,\mathbb R)$, the invertible real matrices, is open in
$\mathrm M(n,\mathbb R)$; multiplication is polynomial in the entries and
inversion is smooth by Cramer's rule.

(b) The positive-determinant subgroup $\mathrm{GL}^{+}(n,\mathbb R)$ is open
in $\mathrm{GL}(n,\mathbb R)$ and has dimension $n^2$.

(c) Every open subgroup of a Lie group, with its inherited structure, is a Lie
group.

(d) $\mathrm{GL}(n,\mathbb C)$ is open in $\mathrm M(n,\mathbb C)$, has real
dimension $2n^2$, and its operations are smooth in the real and imaginary
parts of the entries.

(e) For a finite-dimensional real or complex vector space $V$,
$\mathrm{GL}(V)$ is a Lie group. A basis identifies it with
$\mathrm{GL}(n,\mathbb R)$ or $\mathrm{GL}(n,\mathbb C)$, and a change of
basis acts by the smooth map $A\mapsto BAB^{-1}$.

(f) $\mathbb R$ and $\mathbb R^n$ are additive Lie groups.

(g) $\mathbb C$ and $\mathbb C^n$ are additive Lie groups.

(h) $\mathbb R^{*}$ is a one-dimensional multiplicative Lie group, with open
subgroup $\mathbb R^{+}$.

(i) $\mathbb C^{*}$ is a two-dimensional multiplicative Lie group.

(j) $\mathbb S^1\subset\mathbb C^{*}$ is a Lie group under multiplication; in
angle coordinates its operations are
$(\theta_1,\theta_2)\mapsto\theta_1+\theta_2$ and $\theta\mapsto-\theta$.

(k) Products of Lie groups are Lie groups under componentwise multiplication:
\[
(g_1,\ldots,g_k)(g'_1,\ldots,g'_k)
=(g_1g'_1,\ldots,g_kg'_k).
\]

(l) $\mathbb T^n=(\mathbb S^1)^n$ is an $n$-dimensional abelian Lie group.

(m) A finite or countably infinite group with the discrete topology is a
zero-dimensional Lie group.
\end{example}

\section{Lie Group Homomorphisms}

\begin{definition}
\label{def:7-47-extra-1}
\lean{LieGroupIsomorphism}
\leanok


A Lie group homomorphism $F:G\to H$ is a smooth group homomorphism. It is a Lie group isomorphism when it is a diffeomorphism (then $F^{-1}$ is also a homomorphism).
\end{definition}

\begin{example}
\label{ex:7-4}
\lean{Circle.toUnits}

\mathlibok
\dcref{ch7:7.4}
(a) The inclusion $\mathbb S^1\hookrightarrow\mathbb C^*$ is a Lie group homomorphism.

(b) $\exp:\mathbb R\to\mathbb R^*$, $t\mapsto e^t$, is a Lie group
homomorphism with image $\mathbb R^+$ and restricts to a Lie group isomorphism
$\mathbb R\cong\mathbb R^+$ with inverse $\log$.

(c) $\exp:\mathbb C\to\mathbb C^*$ is a surjective Lie group homomorphism with
kernel $2\pi i\mathbb Z$.

(d) $\varepsilon:\mathbb R\to\mathbb S^1$, $t\mapsto e^{2\pi it}$, is a Lie
group homomorphism with kernel $\mathbb Z$; componentwise,
$\varepsilon^n:\mathbb R^n\to\mathbb T^n$ is a Lie group homomorphism with
kernel $\mathbb Z^n$.

(e) Determinants $\mathrm{GL}(n,\mathbb R)\to\mathbb R^*$ and $\mathrm{GL}(n,\mathbb C)\to\mathbb C^*$ are smooth homomorphisms.

(f) Conjugation $C_g(h)=ghg^{-1}$ is a Lie group automorphism with inverse $C_{g^{-1}}$. A subgroup $H$ is normal exactly when $C_g(H)=H$ for every $g$.
\end{example}

\section{The Universal Covering Group}


\begin{theorem}[Uniqueness of the Universal Covering Group]
\label{thm:7-9}
\lean{exists_lieGroupIsomorphism_of_universal_covering_group}
\uses{prob:7-5}
\leanok


\dcref{ch7:7.9}
For a connected Lie group $G$, if simply connected Lie groups $\widetilde G$ and $\widetilde G'$ admit smooth covering homomorphisms $\pi:\widetilde G\to G$ and $\pi':\widetilde G'\to G$, then there is a Lie group isomorphism $\Phi:\widetilde G\to\widetilde G'$ satisfying $\pi'\circ\Phi=\pi$.
\end{theorem}

\section{Lie Subgroups}

\begin{definition}
\label{def:7-49-extra-1}
\lean{LieSubgroup}
\leanok


A Lie subgroup of a Lie group $G$ is a subgroup with a Lie-group structure whose inclusion in $G$ is an immersion.
\end{definition}

\begin{proposition}
\label{prop:7-11}
\lean{subgroup_has_lieSubgroup_structure_of_isEmbeddedSubmanifold}
\uses{cor:5-30}



\dcref{ch7:7.11}
An embedded submanifold $H\subseteq G$ that is a subgroup is a Lie subgroup.
\end{proposition}

\begin{proof}
\leanok
\sketch Restrict multiplication and inversion on $G$ to $H$. Their images lie in $H$ by the subgroup property, and the embedded-submanifold smoothness criterion (Corollary~\ref{cor:5-30}) makes the restricted maps smooth into $H$.
\end{proof}

\begin{definition}
\label{def:7-49-extra-2}
\lean{subgroupGeneratedBy}
\leanok


For a group $G$ and $S\subseteq G$, the subgroup generated by $S$ is the intersection of all subgroups containing $S$.
\end{definition}

\begin{lemma}[Words in a Generated Subgroup]
\label{exer:7-13}
\lean{subgroup_closure_eq_subgroup_word_set}
\leanok


\dcref{ch7:7.13}
The subgroup generated by $S$ consists exactly of finite products of elements
of $S$ and their inverses.
\end{lemma}

\begin{proposition}
\label{prop:7-14}
\lean{generatedSubgroup_eq_top_of_mem_nhds_one}
\uses{exer:7-13}
\leanok



\dcref{ch7:7.14}
For a Lie group $G$ and identity neighborhood $W$, the subgroup generated by $W$ is open; if $W$ is connected it is connected; and if $G$ is connected it equals $G$.
\end{proposition}

\begin{proof}
\leanok
\sketch Let $H=\langle W\rangle$ and $W_k$ be products of at most $k$ elements of $W\cup W^{-1}$. Since inversion and left translations are diffeomorphisms, induction gives that every $W_k$ and hence $H=\bigcup_kW_k$ is open. If $W$ is connected, then $W\cup W^{-1}$ and all products $W_k$ are connected, so $H$ is connected. An open subgroup is closed; therefore a connected $G$ has $H=G$.
\end{proof}

\begin{definition}
\label{def:7-49-extra-3}
\lean{connectedComponentOfOne_coe}
\leanok


The identity component of a Lie group is the connected component containing its identity.
\end{definition}

\begin{example}
\label{ex:7-18}
\lean{complex_generalLinear_det_surjective}
\uses{ex:7-3}
\leanok



\dcref{ch7:7.18}
(a) $\mathrm{GL}^{+}(n,\mathbb R)$ is an open embedded Lie subgroup of
$\mathrm{GL}(n,\mathbb R)$.

(b) $\mathbb S^1$ is an embedded Lie subgroup of $\mathbb C^*$.

(c) $\mathrm{SL}(n,\mathbb R)
=\ker(\det:\mathrm{GL}(n,\mathbb R)\to\mathbb R^*)$ is properly embedded of
dimension $n^2-1$, since determinant is a surjective submersion.

(d) Replacing each $a+ib$ by
$\left(\begin{smallmatrix}a&-b\\b&a\end{smallmatrix}\right)$ defines an
injective Lie group homomorphism
$\beta:\mathrm{GL}(n,\mathbb C)\to\mathrm{GL}(2n,\mathbb R)$ with properly
embedded image.

(e) $\mathrm{SL}(n,\mathbb C)
=\ker(\det:\mathrm{GL}(n,\mathbb C)\to\mathbb C^*)$ is properly embedded of
codimension $2$ and dimension $2n^2-2$.
\end{example}

\section{Group Actions and Equivariant Maps}

\begin{example}
\label{ex:7-22}
\lean{Example_7_22}
\leanok


\dcref{ch7:7.22}
(a) The trivial action $g\cdot p=p$ of any Lie group $G$ on a manifold $M$ has point orbits and isotropy group $G$.

(b) $\mathrm{GL}(n,\mathbb R)$ acts smoothly on $\mathbb R^n$ by $(A,x)\mapsto Ax$; its orbits are $\{0\}$ and $\mathbb R^n\setminus\{0\}$.

(c) Left translation makes every Lie group act smoothly, freely, and transitively on itself. Restriction gives a smooth free action of any Lie subgroup $H$ on $G$; right translations are analogous.

(d) Conjugation gives the smooth action $g\cdot h=ghg^{-1}$.

(e) For a discrete group $\Gamma$, an action on $M$ is smooth exactly when each map $p\mapsto g\cdot p$ is smooth. In particular, $\mathbb Z^n$ acts smoothly and freely on $\mathbb R^n$ by translation.
\end{example}

\begin{definition}
\label{def:7-50-extra-2}
\lean{coveringAutomorphismGroup}
\leanok


For a covering map $\pi:E\to M$, a covering (deck) transformation is a
homeomorphism $\varphi:E\to E$ with $\pi\circ\varphi=\pi$. These transformations
form the group $\operatorname{Aut}_\pi(E)$ under composition and act on $E$ by
$\varphi\cdot q=\varphi(q)$. The action is transitive on each fiber exactly
when $\pi_*(\pi_1(E,q))\triangleleft\pi_1(M,\pi(q))$ for every $q$ (the
normal-covering condition).
\end{definition}

\begin{proposition}
\label{prop:7-23}
\lean{coveringAutomorphismGroup_isFree}
\leanok


\dcref{ch7:7.23}
If $\pi:E\to M$ is a smooth covering of smooth manifolds with or without
boundary, then $\operatorname{Aut}_\pi(E)$ with the discrete topology is a
zero-dimensional Lie group acting smoothly and freely on $E$.
\end{proposition}

\begin{proof}
\leanok
\sketch A deck transformation fixing a point agrees there with the identity lift of $\pi$, so uniqueness of lifts makes it the identity; the action is free. For an evenly covered neighborhood $U$ of $\pi(q)$, second countability makes the fiber $\pi^{-1}(\pi(q))$ countable. The injective evaluation map $\varphi\mapsto\varphi(q)$ embeds $\operatorname{Aut}_\pi(E)$ in this fiber, so the group is countable and hence a discrete zero-dimensional Lie group. Smoothness of the action follows from Theorem~4.29.
\end{proof}

\begin{definition}
\label{def:7-50-extra-3}
\lean{MulActionHom}

\mathlibok
Given smooth left (respectively right) $G$-actions on $M$ and $N$, a map $F:M\to N$ is equivariant when $F(g\cdot p)=g\cdot F(p)$ (respectively $F(p\cdot g)=F(p)\cdot g$) for every $g$. Equivalently, it intertwines the two actions.
\end{definition}

\begin{definition}
\label{def:7-50-extra-4}
\lean{orbit_map}
\leanok


For a smooth left action $\theta:G\times M\to M$ and $p\in M$, the orbit map is $\theta^{(p)}:G\to M$, $\theta^{(p)}(g)=g\cdot p$. Its image is $G\cdot p$ and $(\theta^{(p)})^{-1}(p)$ is the isotropy group $G_p$.
\end{definition}

\begin{example}[The Orthogonal Group]
\label{ex:7-27}
\lean{isCompact_orthogonalGroup}
\leanok


\dcref{ch7:7.27}
The orthogonal group $\mathrm O(n)=\{A\in\mathrm{GL}(n,\mathbb R):A^TA=I_n\}$ consists of matrices preserving the Euclidean inner product. The map $\Phi(A)=A^TA$ is equivariant for right multiplication on $\mathrm{GL}(n,\mathbb R)$ and $X\cdot B=B^TXB$ on matrices, so $\mathrm O(n)=\Phi^{-1}(I_n)$ is a properly embedded Lie subgroup. At $I_n$, $d\Phi_{I_n}(B)=B^T+B$, whose image is the space of symmetric matrices of dimension $n(n+1)/2$; hence $\dim\mathrm O(n)=n(n-1)/2$. Every $A\in\mathrm O(n)$ has orthonormal columns and Frobenius norm $\lvert A\rvert=\sqrt n$, so the subgroup is compact as closed and bounded.
\end{example}

\begin{example}[The Unitary Group]
\label{ex:7-29}
\lean{Matrix.conjTranspose_mul}

\mathlibok
\dcref{ch7:7.29}
For a complex matrix $A$, write $A^*=\overline A^T$, so $(AB)^*=B^*A^*$. With Hermitian product $z\cdot w=\sum_i z^i\overline{w^i}$, the unitary group $\mathrm U(n)=\{A\in\mathrm{GL}(n,\mathbb C):A^*A=I_n\}$ consists of matrices with orthonormal columns and is a properly embedded Lie subgroup of dimension $n^2$.
\end{example}

\begin{example}[The Special Unitary Group]
\label{ex:7-30}
\lean{Matrix.specialUnitaryGroup}

\mathlibok
\dcref{ch7:7.30}
$\mathrm{SU}(n)=\mathrm U(n)\cap\mathrm{SL}(n,\mathbb C)$ is a properly embedded Lie subgroup of $\mathrm U(n)$ of dimension $n^2-1$, and hence is embedded in $\mathrm{GL}(n,\mathbb C)$.
\end{example}

\section{Semidirect Products}

\begin{definition}
\label{def:7-51-extra-1}
\lean{semidirectProductGroup_mul_eq}
\leanok


Let $H,N$ be Lie groups and let $\theta:H\times N\to N$ be a smooth action by automorphisms. The semidirect product $N\rtimes_\theta H$ is the manifold $N\times H$ with
\[
(n,h)(n',h')=\bigl(n\theta_h(n'),hh'\bigr).
\]
\end{definition}

\begin{proposition}[Semidirect Products Are Lie Groups]
\label{exer:7-31}
\lean{semidirectProductLieGroup}
\leanok


\dcref{ch7:7.31}
The multiplication in Definition~\ref{def:7-51-extra-1} defines a Lie group
with identity $(e,e)$ and inverse
$(n,h)^{-1}=(\theta_{h^{-1}}(n^{-1}),h^{-1})$.
\end{proposition}

\begin{example}[The Euclidean Group]
\label{ex:7-32}
\lean{euclidean_group_affine_equiv}
\leanok


\dcref{ch7:7.32}
The Euclidean group is $\mathrm E(n)=\mathbb R^n\rtimes\mathrm O(n)$ with
\[
(b,A)(b',A')=(b+Ab',AA'),\qquad (b,A)\cdot x=b+Ax.
\]
This action preserves lines, distances, and angles.
\end{example}

\begin{proposition}[Properties of Semidirect Products]
\label{prop:7-33}
\lean{proposition_7_33_tildeN_mul_tildeH}
\leanok


\dcref{ch7:7.33}
For $G=N\rtimes_\theta H$ as in Definition~\ref{def:7-51-extra-1},
$\widetilde N=N\times\{e\}$ and $\widetilde H=\{e\}\times H$ are closed Lie
subgroups isomorphic to $N$ and $H$; $\widetilde N$ is normal in $G$,
$\widetilde N\cap\widetilde H=\{(e,e)\}$, and
$\widetilde N\widetilde H=G$.
\end{proposition}


\section{Representations}

\begin{definition}
\label{def:7-52-extra-1}
\lean{LieGroupRepresentation.isFaithful_iff_injective}
\leanok


A finite-dimensional representation of a Lie group $G$ on a real or complex vector space $V$ is a Lie group homomorphism $\rho:G\to\mathrm{GL}(V)$. It is faithful when injective. A faithful representation identifies $G$ with a Lie subgroup of some $\mathrm{GL}(n,\mathbb R)$ or $\mathrm{GL}(n,\mathbb C)$; not every Lie group has one (the universal cover of $\mathrm{SL}(2,\mathbb R)$ is a stated counterexample).
\end{definition}

\begin{definition}
\label{def:7-52-extra-2}
\lean{IsLinearAction}
\leanok


A smooth action of $G$ on a finite-dimensional vector space $V$ is linear if each map $x\mapsto g\cdot x$ is linear. A representation $\rho$ induces the linear action $g\cdot x=\rho(g)x$.
\end{definition}

\begin{proposition}
\label{prop:7-37}
\lean{isLinearAction_iff_exists_lieGroupRepresentation}
\leanok


\dcref{ch7:7.37}
A smooth left action of a Lie group $G$ on a finite-dimensional vector space $V$ is linear if and only if $g\cdot x=\rho(g)x$ for a representation $\rho:G\to\mathrm{GL}(V)$.
\end{proposition}

\begin{proof}
\leanok
\sketch A linear action defines invertible linear maps $\rho(g)$ and the action law gives $\rho(g_1g_2)=\rho(g_1)\rho(g_2)$. In a basis $(E_j)$ with coordinate projections $\pi^i$, each matrix coefficient is $\rho^i_j(g)=\pi^i(g\cdot E_j)$, hence smooth; therefore $\rho$ is a Lie group homomorphism. The converse is immediate.
\end{proof}

\begin{exercise}
\label{exer:7-2}
\lean{lieGroup_of_contMDiff_mul_inv}
\leanok
\uses{prop:7-1}
\dcref{ch7:7.2}
\begin{itemize}
\item Exercise 7.2. Prove Proposition 7.1.
\end{itemize}
\end{exercise}

\begin{exercise}
\label{exer:7-34}
\lean{proposition_7_33_tildeN_mul_tildeH}
\leanok
\uses{prop:7-33}
\dcref{ch7:7.34}
\begin{itemize}
\item Exercise 7.34. Prove the preceding proposition.
\end{itemize}
\end{exercise}

\begin{exercise}
\label{exer:7-8}
\lean{mul_inv_cancel}
\mathlibok
\dcref{ch7:7.8}
Complete the proof of the preceding theorem by showing that ${x}^{-1}x = x{x}^{-1} = \widetilde{e}.$
\end{exercise}
